\part{Volume 2}

\chapter{Introduction to classical field theory}

I wanted to begin some light reading of concepts covered in volume 2, such as $L_\infty $-algebras, elliptic moduli problems, etc. Since I'm only reading for a few days, I am not making a new project for this, and am keeping it in this file. We will mostly be following part 1 of \cite{cg2}, as well as their appendicies. We will be using significant supplimental material since they provide only a brief overview of many important topics.

We start by following chapter 2 of \cite{cg2} for motivation. However, since the current focus is not on understanding the mathematical physics, but rather the mathematical tools, we will keep this chapter very brief.

We will be working perturbatively throughout this book. Our goal is to describe the derived space of solutions to the Euler-Lagrange equations. Since these are usually a system of PDEs on a manifold $M$, and those PDEs are usually elliptic in a particular sense, the space we want to describe will be called a \emph{formal elliptic moduli problem}. Since one can solve the Euler-Lagrange equations on any open subset $U\subseteq M$, the correct object to use is a sheaf $\mathcal{EL} $ of formal elliptic moduli problems on $M$.

We then think of observables as some kind of function on $\mathcal{E} \mathcal{L} $, so that
\[
	\Obs^{cl} (U) = \mathcal{O} (\mathcal{E} \mathcal{L} (U))
.\]
We will be precise about what this means later. The factorization algebra structure arises from this idea. In finite dimensions, a formula moduli problem that arises as the derived critical locus of a function is equipped with a symplectic form of cohomological degree $-1$. This symplectic form is quite useful, and in fact we later \emph{define} a classical field theory to be a formal elliptic moduli problem equipped with a symplectic form of cohomological degree $-1$. We can then show how to construct such a theory from a local action functional satisfying certain non-degeneracy properties. Finally, we look at a particular construction on an elliptic moduli problem $\mathcal{F} $ similar to that of the cotangent bundle. This is motivated by symplectic geometry, and turns out to be physically interesting.

In finite dimensions, the symplectic form for a formal moduli problem $\mathcal{F} $ yields a Poisson bracket of degree 1 for $\mathcal{O} (\mathcal{F} )$, hence $\mathcal{O} (\mathcal{F} )$ is a $P_0$-algebra. We show that something similar happens in infinite dimensions: any classical field theory has a $P_0$-structure on its factorization algebra of observables. This is the main payoff of ellipticity of the Euler-Lagrange equations.

\chapter{$L_\infty $ algebras and moduli problems}

\section{Lie algebras and $L_\infty $ algebras}

Given a dg Lie algebra $\mathfrak{g} $, we have extensively used the Chevalley-Eilenberg chains and cochains to construct many important factorization algebras. It is necessary to make a homotopic generalization of the notion of a Lie algebra for interacting theories, and the category $\Ch_{\mathbb{K} } $ (here $\mathbb{K}$ is a field of characteristic 0) is a great place to work, since it is an $\infty $-category via the dg nerve.

\begin{definition}\label{dfn:10-1-1}
	An $L_\infty $ algebra over a ring $R$ is a graded projective $R$-module $\mathfrak{g} $ equipped with multilinear maps
	\[
		\ell_n : \mathfrak{g} \otimes_R \cdots \otimes _R \mathfrak{g}  \to \mathfrak{g} [n-2]
	\]
	for all $n\geq 1$, referred to as \emph{brackets}, satisfying the following properties.
	\begin{enumerate}
		\item Each bracket $\ell _n$ is graded-alternating, meaning
			\[
				\ell _n(x_1, \ldots ,x_i,x_{i+1} ,\ldots ,x_n) = -(-1)^{\left| x_i \right| \left| x_{i+1}  \right| } \ell _n(x_1, \ldots , x_{i+1} ,x_i, \ldots ,x_n)
			.\]
		\item Each bracket $\ell _n$ satisfies the $n$-Jacobi rule: for all tuples $(x_1, \ldots , x_n)\in \mathfrak{g} $,
			\[
				0=\sum_{k=1}^{n} (-1)^k \sum_{\substack{i_1 < \cdots < i_k \\ j_{k+1} < \cdots < j_n \\ \left\{ i_1, \ldots , j_n \right\} =\left\{ 1, \ldots ,n \right\} }} \varepsilon_{i_1, \ldots ,j_n} \ell _{n-k+1} (\ell _k(x_{i_1} , \ldots ,x_{i_k} ), x_{j_{k+1} } , \ldots , x_{j_n} )
			.\]
			Here, $\varepsilon $ yields the proper sign for the alternating Koszul sign rule $ab=-(-1)^{\left| a \right| \left| b \right| } ba$.
	\end{enumerate}	
\end{definition}

Part (2) is a tad confusing so I shall explain it. To be more explicit about the actual equation itself, I will write down some terms from the sum. We will let $x_1,x_2,x_3\in \mathfrak{g} $, and work with the 3-Jacobi rule. First, we set $k=1$. Then we have $\left\{ i_1, j_2, j_3 \right\} =\left\{ 1,2,3 \right\} $. If we set $i_1=1$, $j_2=2$, and $j_3=3$, then we simply have
\[
	(-1)^k \varepsilon \ell_3(\ell _1(x_1),x_2,x_3)
.\]
Since there are no crossings, we have $\varepsilon =-(-1)^1 = 1$. Thus this term evaluates to simply
\[
	\ell _3(\ell _1(x_1),x_2,x_3)
.\]
If we were to switch $x_1$ and $x_2$ by setting $i_1 = 2$ and $j_2 = 1$, then we would have $\varepsilon =(-1)^{\left| x_1 \right| \left| x_2 \right|}$, yielding the term
\[
	(-1)^k {\varepsilon } \ell _3(\ell _1(x_2),x_1,x_3) = -(-1)^{\left| x_1 \right| \left| x_2 \right| } \ell _3(\ell _1(x_2),x_1,x_3)
.\]
From here, it should be clear how the rest of the terms are calculated.

We can use this to compute low dimensional examples of the brackets. For example, the 1-Jacobi rule simply says
\[
	0 = \sum_{k=1}^{1} (-1)^k\sum_{i_1 = 1} \varepsilon_{i_1} \ell_{1 - k + 1} (\ell_k(x_{i_1})) = -(-1)^{1} \ell _1(\ell _1(x_1))
.\]
Since there are no crossings, $\varepsilon =1$, so this evaluates to
\[
	0 = \ell_1(\ell _1(x_1))
,\]
or more simply, $\ell _1\circ \ell _1=0$. Thus, $\ell _1$ is a differential. Let's temporarily write $\dd =\ell _1$ and $[-,-]=\ell _2$. We now consider the 2-Jacobi rule:
\[
	0=\sum_{k=1}^{2} (-1)^k \sum_{\substack{i_1 < \cdots < i_k \\ j_{k+1} < \cdots < j_n \\ \left\{ i_1, \ldots , j_n \right\} =\left\{ 1, \ldots ,n \right\} }} \varepsilon_{i_1, \ldots ,j_n}  \ell _{n-k+1} (\ell _k(x_{i_1} , \ldots ,x_{i_k} ), x_{j_{k+1} } , \ldots , x_{j_n} )
.\]
There are two possible values for $k$, so we work with each one individually. First, consider the terms where $k=1$. Let $x_1,x_2\in \mathfrak{g} $. The list of all possible permutations, and the corresponding sign value, is
\begin{align*}
	(x_1,x_2) && \varepsilon =-(-1)^1=1\\
	(x_2,x_1) && \varepsilon =-(-1)^{\left| x_1 \right| \left| x_2 \right| } 
.\end{align*}
The $k=1$ component of the summand ranges over $i_1=1$ and $i_1=2$, hence it ranges over both permutations. Since $(-1)^k=-1$, the $k=1$ component is simply
\[
	-[\dd x_1,x_2] + (-1)^{\left| x_1 \right| \left| x_2 \right| } [\dd x_2,x_1]
.\]
The $k=2$ component of the summand forces $i_1=1$ and $i_2=2$ due to the reqirement that $i_1 < i_2$. Thus, it is simply
\[
	(-1)^2(- (-1)^1) \ell_{2-2+1} (\ell _2(x_1,x_2)) = \dd [x_1,x_2]
.\]
Putting it together, we have
\[
	0=\dd [x_1,x_2] - [\dd x_1,x_2] + (-1)^{\left| x_1 \right| \left| x_2 \right| } [\dd x_2,x_1]
.\]
We can write this in a more conventional form. First, we move $\dd[x_1,x_2]$ over to the other side to get
\[
	-\dd [x_1,x_2] = -[\dd x_1,x_2] + (-1)^{\left| x_1 \right| \left| x_2 \right| } [\dd x_2,x_1]
.\]
We then multiply through by $-1$ to obtain
\[
	\dd [x_1,x_2] = [\dd x_1,x_2] - (-1)^{\left| x_1 \right| \left| x_2 \right| } [\dd x_2,x_1] 
.\]
Recall that $\ell _2$ is graded alternating, meaning $[x,y] = -(-1)^{\left| x \right| \left| y \right| } [y,x]$. Since $\left| \dd x_2 \right| =\left| x_2 \right| +1$, we have that
\[
	-(-1)^{\left| x_1 \right| \left| x_2 \right| } =-(-1)^{\left| x_1 \right| \left| x_2 \right| } (-(-1)^{\left| x_1 \right| \left| x_2 \right| +\left| x_1 \right| } )[x_1,\dd x_2]=(-1)^{2\left| x_1 \right| \left| x_2 \right| +\left| x_1 \right| } [x_1,\dd x_2]=(-1)^{\left| x_1 \right| } [x_1,\dd x_2]
.\]
Thus the 2-Jacobi rule reads
\[
	\dd [x_1,x_2] = [\dd x_1,x_2] + (-1)^{\left| x_1 \right| } [x_1,\dd x_2]
.\]
This is precisely the graded Leibniz rule. Additionally, note that due to the definition of the differential on the graded tensor product $\mathfrak{g} \otimes_R \mathfrak{g} $, this is also the requirement that $\ell _2$ be a cochain map $\mathfrak{g} \otimes _R\mathfrak{g} \to \mathfrak{g} $.

Now we consider the 3-Jacobi rule. Let $x_1,x_2,x_3\in \mathfrak{g} $. Then the possible permutations, and their corresponding $\varepsilon $ value are
\begin{align*}
	(x_1,x_2,x_3)&& \varepsilon =-(-1)^{1}=1\\
	(x_1,x_3,x_2) && \varepsilon =-(-1)^{\left| x_2 \right| \left| x_3 \right| } \\
	(x_2,x_1,x_3) && \varepsilon =-(-1)^{\left| x_1 \right| \left| x_2 \right|}\\
	(x_2,x_3,x_1) && \varepsilon =(-1)^{\left| x_1 \right| \left| x_2 \right|}(-1)^{\left| x_1 \right| \left| x_3 \right| } \\
	(x_3,x_2,x_1) && \varepsilon =-(-1)^{\left| x_1 \right| \left| x_2 \right| }(-1)^{\left| x_1 \right| \left| x_3 \right| } (-1)^{\left| x_2 \right| \left| x_3 \right| } \\
	(x_3,x_1,x_2) && \varepsilon =(-1)^{\left| x_1 \right| \left| x_3 \right| } (-1)^{\left| x_2 \right| \left| x_3 \right| } 
.\end{align*}
Note that not all of these are possible, since we must have $i_1 < i_2 < i_3$ and likewise for the $j$'s. In fact, the only possible permutations are
\[
	(i_1,i_2,i_3) = (i_1,i_2,j_3) = (i_1,j_2,j_3) = (x_1,x_2,x_3),\qquad (i_1,i_2,j_3) = (x_1,x_3,x_2),
\]
\[
	(i_1,j_2,j_3) = (x_2,x_1,x_3),\qquad (i_1,i_2,j_3) = (x_2,x_3,x_1),\qquad (i_1,j_2,j_3) = (x_3,x_1,x_2)
.\]
Essentially, all permutations except $(x_3,x_2,x_1)$ are possible. We once again work termwise under $k$. Start with $k=3$. Then we must have 3 $i$'s, meaning there is only one possible permutation, $(i_1,i_2,i_3) = (x_1,x_2,x_3)$. Thus the $k=3$ term is
\[
	(-1)^3(-(-1)^1) \ell_{3-3+1} (\ell _3(x_1,x_2,x_3)) = -\dd \ell _3(x_1,x_2,x_3)
.\]
The bracket for the $k=2$ terms is $\ell _{3-2+1} \circ \ell_2 = \ell_2(\ell _2,-)$. Thus plugging in the terms with two $i$'s yields the $k=2$ terms:
\[
	[[x_1,x_2],x_3] - (-1)^{\left| x_2 \right| \left| x_3 \right| } [[x_1,x_3],x_2] + (-1)^{\left| x_1 \right| (\left| x_2 \right| +\left| x_3 \right|)} [[x_2,x_3],x_1]
.\]
We will later be able to rearrange the brackets to get rid of these signs. For now, we compute the $k=1$ terms. These all pick up a factor of $-1$ from $(-1)^k=-1$, which cancels the leading minus sign in all $\varepsilon $'s except for the $(x_1,x_2,x_3)$ permutation, where $\varepsilon =1$. Moreover, $\ell _{3-k+1} =\ell _3$, so the operator is simply $\ell _3(\dd ,-,-)$. Taking all permutations with two $j$'s yields
\[
	-\ell _3(\dd x_1,x_2,x_3) + (-1)^{\left| x_1 \right| \left| x_2 \right| } \ell _3(\dd x_2,x_1,x_3) - (-1)^{\left| x_3 \right| (\left| x_1 \right| +\left| x_2 \right| )} \ell _3(\dd x_3,x_1,x_2)
.\]
We can now write the full 3-Jacobi rule:
\begin{align*}
	0&=-\dd \ell _3(x_1,x_2,x_3) + [[x_1,x_2],x_3] - (-1)^{\left| x_2 \right| \left| x_3 \right| } [[x_1,x_3],x_2] + (-1)^{\left| x_1 \right| (\left| x_2 \right| +\left| x_3 \right| )} [[x_2,x_3],x_1]\\
	 &\,\,\,\,\,\,\,-\ell _3(\dd x_1,x_2,x_3) + (-1)^{\left| x_1 \right| \left| x_2 \right| } \ell _3(\dd x_2,x_1,x_3) - (-1)^{\left| x_3 \right| (\left| x_1 \right| +\left| x_2 \right| )} \ell _3(\dd x_3,x_1,x_2)
.\end{align*}
We will isolate the terms with $\ell _2$ brackets from the terms with $\ell _3$ brackets to begin putting this in a more conventional form. We also use $[x_1,x_3] = -(-1)^{\left| x_1 \right| \left| x_3 \right| } [x_3,x_1]$. We rewrite
\begin{align*}
	&[[x_1,x_2],x_3] + (-1)^{\left| x_3 \right|(\left| x_1 \right| +\left| x_2 \right| ) } [[x_3,x_1],x_2] + (-1)^{\left| x_1 \right| (\left| x_2 \right| +\left| x_3 \right| )} [[x_2,x_3],x_1]\\
	&=\dd \ell _3(x_1,x_2,x_3) + \ell _3(\dd x_1,x_2,x_3) - (-1)^{\left| x_1 \right| \left| x_2 \right| } \ell _3(\dd x_2,x_1,x_3) + (-1)^{\left| x_3 \right| (\left| x_1 \right| +\left| x_2 \right| )} \ell _3(\dd x_3,x_1,x_2)
.\end{align*}
We will hereafter refer to this expression as the 3-Jacobi rule, and will make specific reference to its left hand side and right hand side.

We claim that this equation encodes the usual graded Jacobi identity up to homotopy, and this homotopy is encoded in the bracket $\ell _3$. Thus, we should first check that when $\ell _3=0$, we recover the usual graded Jacobi identity
\[
	(-1)^{\left| x \right| \left| z \right| } [x,[y,z]]+(-1)^{\left| y \right| \left| x \right| } [y,[z,x]]+(-1)^{\left| z \right| \left| y \right| } [z,[x,y]]=0
.\]
Setting $\ell _3=0$ yields
\[
	[[x_1,x_2],x_3] + (-1)^{\left| x_3 \right|(\left| x_1 \right| +\left| x_2 \right| ) } [[x_3,x_1],x_2] + (-1)^{\left| x_1 \right| (\left| x_2 \right| +\left| x_3 \right| )} [[x_2,x_3],x_1]=0
.\]
Note that $\left| [x,y] \right| =\left| x \right| +\left| y \right| $ since $\ell _2$ has cohomological degree 0. Thus,
\[
	[x,[y,z]] = -(-1)^{\left| x \right| (\left| y \right| +\left| z \right| )} [[y,z],x]
.\]
We can apply this to the first and third term of our equation, allowing us to write
\[
	-(-1)^{\left| x_3 \right| (\left| x_1 \right| +\left| x_2 \right| )} [x_3,[x_1,x_2]] + (-1)^{\left| x_3 \right| (\left| x_1 \right| +\left| x_2 \right| )} [[x_3,x_1],x_2] -[x_1,[x_2,x_3]]=0
.\]
We did not apply this to the second term, because we would like to treat that term separately. We have
\[
	(-1)^{\left| x_3 \right| (\left| x_1 \right| +\left| x_2 \right| )} [[x_3,x_1],x_2] = -(-1)^{\left| x_3 \right| (\left| x_1 \right| +\left| x_2 \right| )+\left| x_2 \right| (\left| x_3 \right| +\left| x_1 \right| )} [x_2,[x_3,x_1]] = -(-1)^{\left| x_1 \right( \left| x_3 \right|+ \left| x_2 \right|)} [x_2,[x_3,x_1]]
.\]
The last equality is given by expanding out the exponent and removing the factor of $2\left| x_2 \right| \left| x_3 \right| $. Substituting this in and distributing in all the exponents, we have
\[
	-(-1)^{\left| x_3 \right| \left| x_1 \right| +\left| x_3 \right| \left| x_2 \right| } [x_3,[x_1,x_2]] -(-1)^{\left| x_1 \right( \left| x_3 \right|+ \left| x_2 \right|)} [x_2,[x_3,x_1]] -[x_1,[x_2,x_3]]=0
.\]
Note that both terms with an exponent contain an $\left| x_1 \right| \left| x_3 \right| $ term. We can now simply multiply through by $-(-1)^{\left| x_3 \right| \left| x_1 \right| } $ to obtain
\[
	(-1)^{\left| x_2 \right| \left| x_3 \right| } [x_3,[x_1,x_2]]+(-1)^{\left| x_1 \right| \left| x_2 \right| } [x_2,[x_3,x_1]]+(-1)^{\left| x_1 \right| \left| x_3 \right| } [x_2,[x_3,x_1]]=0
.\]
This is precisely the graded Jacobi identity. Thus, if $\mathfrak{g} $ is an $L_\infty $ algebra with $\ell_n=0$ for all $n\geq 3$, then $\mathfrak{g} $ is precisely a differential graded Lie $R$-algebra.

Now, we will consider the case where $\ell _3 \neq 0$. We claim that the datum of $\ell _3$ allows us to pass to cohomology, and to have the graded Jacobi identity hold in cohomology. First, we note that if $x=\dd \eta $ is exact and $y$ is closed, then by the graded Leibniz rule,
\[
	[x,y]=[\dd \eta ,y] = \dd [\eta ,y] - (-1)^{\left| \eta  \right| } [\eta ,\dd y]=\dd [\eta ,y] - 0 = \dd [\eta ,y]
.\]
Thus, $[x,y]$ is exact. This is the standard argument that a dgla's cohomology inherits a well-defined bracket: given any exact form, the bracket with that exact form is necessarily killed.

We can apply similar logic to the 3-Jacobi rule. Let $x_1,x_2,x_3$ now be closed elements, such that they define a cohomology class. The right hand side of the 3-Jacobi rule reads
\[
	\dd \ell _3(x_1,x_2,x_3) + \ell _3(\dd x_1,x_2,x_3) - (-1)^{\left| x_1 \right| \left| x_2 \right| } \ell _3(\dd x_2,x_1,x_3) + (-1)^{\left| x_3 \right| (\left| x_1 \right| +\left| x_2 \right| )} \ell _3(\dd x_3,x_1,x_2)
.\]
Since each $x_i$ is closed, $\dd x_i=0$ for all $i$, so the latter three $\ell _3$-brackets necessarily vanish, leaving just
\[
	\dd \ell _3(x_1,x_2,x_3)
.\]
Note that this is an exact element. Thus, when evaluating on closed elements, we have
\[
	[[x_1,x_2],x_3] + (-1)^{\left| x_3 \right|(\left| x_1 \right| +\left| x_2 \right| ) } [[x_3,x_1],x_2] + (-1)^{\left| x_1 \right| (\left| x_2 \right| +\left| x_3 \right| )} [[x_2,x_3],x_1] = \dd \ell _3(x_1,x_2,x_3)
.\]
Alternatively, using the more conventional graded Jacobi identity (and recalling that to derive it, we had to multiply through by $(-1)^{\left| x_1 \right| \left| x_3 \right| } $), we can write
\[
	(-1)^{\left| x_1 \right| \left| x_3 \right| } [x_1,[x_2,x_3]]+(-1)^{\left| x_1 \right| \left| x_2 \right| } [x_2,[x_3,x_1]]+(-1)^{\left| x_2 \right| \left| x_3 \right| } [x_3,[x_1,x_2]]=\dd \ell _3(x_1,x_2,x_3)
.\]
Regardless of how we write it, the point is that the left hand side is an exact element, and thus vanishes in cohomology. Therefore, the Jacobi identity is satisfied in cohomology.

Finally, we can prove that $\ell_3$ is a homotopy which guarentees this relation. This is a trivial application of basic results, which we now recall. First, a cochain map $f\simeq 0$ if and only if there is a homotopy $h$ with $\mathrm{D} (h)=f$, for $\mathrm{D} $ the differential on the internal hom $\underline{\Ch}_R(\dom f, \cod f)$. In this case, we use $\underline{\Ch}_R(\mathfrak{g} ^{\otimes 3} ,\mathfrak{g} )$. Next, recall that since a homotopy $h$ has cohomological degree $-1$, we have $\mathrm{D} (h) = \dd_{\mathfrak{g} } \circ h + h \circ  \mathrm{d}_{\mathfrak{g}^{\otimes 3} } $. Finally, recall that given differential graded $R$-modules $M,N$, the differential is given by
\[
	\dd _{M\otimes_RN}(m\otimes n) = \dd _M(m)\otimes n + (-1)^{\left| m \right| } m\otimes \dd_N(n)
.\]
Applying this here, we have
\[
	\mathrm{D} (\ell _3)(x_1,x_2,x_3) = \dd(\ell _3(x_1,x_2,x_3)) + \ell _3(\dd x_1 \otimes x_2 \otimes x_3 + (-1)^{\left| x_1 \right| } x_1 \otimes \dd x_2 \otimes x_3 +  (-1)^{\left| x_1 \right| +\left| x_2 \right| } x_1\otimes x_2\otimes \dd x_3)
.\]
Since $\ell _3$ is 3-multilinear, and thus linear in triple tensors, this can be rewritten as
\[
	\mathrm{D} (\ell _3)(x_1,x_2,x_3) = \dd \ell _3(x_1,x_2,x_3) + \ell _3(\dd x_1,x_2,x_3) + (-1)^{\left| x_1 \right| } \ell _3(x_1,\dd x_2,x_3) + (-1)^{\left| x_1 \right| +\left| x_2 \right| } \ell _3(x_1,x_2,\dd x_3)
.\]
We claim that we can exactly recover the right hand side of the 3-Jacobi rule from this expression. We have already recovered the first two terms, so we simply need to match the signs on the remaining two. For the first one, we have
\[
	(-1)^{\left| x_1 \right| } \ell _3(x_1,\dd x_2,x_3) = -(-1)^{\left| x_1 \right| +\left| x_1 \right| +\left| x_1 \right| \left| x_2 \right| } \ell _3(\dd x_2,x_1,x_3) = -(-1)^{\left| x_1 \right| \left| x_2 \right| } \ell _3(\dd x_2,x_1,x_3),
\]
which is exactly the term that shows up in the 3-Jacobi rule. Finally, the last remaining term
\[
	(-1)^{\left| x_1 \right| +\left| x_2 \right| } \ell _3(x_1,x_2,\dd x_3) = (-1)^{\left| x_1 \right| +\left| x_2 \right| +\left| x_2 \right| \left| x_3 \right| +\left| x_2 \right| +\left| x_1 \right| \left| x_3 \right| +\left| x_1 \right| } \ell _3(\dd x_3,x_1,x_2)=(-1)^{\left| x_3 \right| (\left| x_1 \right| +\left| x_2 \right| )} \ell _3(\dd x_3,x_1,x_2),
\]
which is the exact final term in the 3-Jacobi rule. Therefore, $\ell _3$ is a \emph{cochain homotopy} expressing the fact that the usual graded Jacobi identity is null-homotopic. We have now finally done good on our motivation of $L_\infty $ algebras as homotopy Lie algebras.

Not all $L_\infty $ algebras have all brackets nonvanishing. In fact, many examples only have a finite number of nonzero brackets. For example, the \emph{String Lie 2-algebra} $\operatorname{string} (n)$ is defined as the graded vector space $\mathfrak{so} (n)\oplus \mathbb{R} \beta [1]$, equipped with two nontrivial brackets
\[
	\ell _2(x,y) = 
	\begin{cases}
		[x,y],&x,y,\in \mathfrak{so} (n)\\
		0,& x=\beta 
	\end{cases}
;\]
\[
	\ell _3(x,y,z) =\mu (x,y,z)\beta ,\qquad x,y,z\in \mathfrak{so} (n),
\]
where $\mu =\left\langle -,[-,-] \right\rangle $ the canonical 3-cocycle arising from the Killing form.

The Chevalley-Eilenberg chains and cochains generalize to $L_\infty $ algebras, and are a particularly important definition.

\begin{definition}\label{dfn:10-1-2}
	Let $\mathfrak{g} $ be an $L_\infty $ algebra over $R$. Then the \emph{Chevalley-Eilenberg chains} for $\mathfrak{g} $, denoted $C_{\bullet} (\mathfrak{g} )$, is the dg cocommutative $R$-coalgebra
	\[
		C_{\bullet} (\mathfrak{g} )\coloneqq \Sym_R(\mathfrak{g} [1])
	.\]
	On generators, the differential $d_n : \Sym ^n_R(\mathfrak{g} [1]) \to \mathfrak{g} [1]$ is given by the bracket $\ell _n$, which we extend to a map on the full symmetric algebra by acting as a coderivation.
\end{definition}

As always, although we call this the Chevalley-Eilenberg \emph{chains}, we view it as a cochain complex to match our conventions. The comultiplication on this coalgebra is given by
\[
	\Delta : C_{\bullet} (\mathfrak{g} ) \to C_{\bullet} (\mathfrak{g} )\otimes_R C_{\bullet} (\mathfrak{g} )
\]
\[
	\Delta (x_1\cdots x_n) = \sum_{\sigma \in S_n} \sum_{1\leq k\leq n-1} (x_{\sigma (1)} \cdots x_{\sigma (k)} )\otimes (x_{\sigma (k+1)}\cdots x_{\sigma (n)} )
.\]
In other words, comultiplication sums over all ways to break a monomial into smaller monomials.

Note that we have not defined a morphism of $L_\infty $ algebras. It would be expected that such a morphism is simply a graded vector space map which respects all brackets. However, it's not clear if such a definition is homotopy-coherent. However, the Chevalley-Eilenberg chains provide a concise way of expressing this information. Note that $C_{\bullet} (\mathfrak{g} )$ encodes all the data of the brackets in its differential. Moreover, it can be shown (tediously) that $\dd ^2=0$ is the datum of the $n$-Jacobi rules. Thus,
\begin{definition}\label{dfn:10-1-3}
	Let $\mathfrak{g} ,\mathfrak{h} $ be $L_\infty $ algebras over $R$. Then a morphism $f : \mathfrak{g}  \to \mathfrak{h} $ is defined as a morphism of differential graded commutative $R$-algebras $f : C_{\bullet} (\mathfrak{g} ) \to C_{\bullet} (\mathfrak{h} )$ on the Chevalley-Eilenberg chains. 
\end{definition}

This is a homotopy-coherent definition. Unpacking the definition shows that such a map decomposes as maps $\Sym ^n_R(\mathfrak{g} [1]) \to \mathfrak{h} $ for each $n$. This is easy to see: the map $\ell _n$ sends any monomial to a single element of $\mathfrak{g} $, so we obtain $f\circ \ell _n = \ell _1 \circ f$.

\begin{definition}\label{dfn:10-1-4}
	Let $\mathfrak{g} $ be an $L_\infty $ algebra over $R$. Then the Chevalley-Eilenberg cochains $C^{\bullet} (\mathfrak{g} )$ are defined as the differential graded commutative $R$-algebra
	\[
		\widehat{\Sym }_R(\mathfrak{g} [1]^\vee )\coloneqq \prod_{n=0}^{\infty } \left( (\mathfrak{g}[1]^{\vee } )^{\otimes n}  \right)_{S_n} 
	.\]
	The differential is given by the derivation $\dd $ whose components $\mathfrak{g} [1]^{\vee } \to \Sym ^n(\mathfrak{g} [1]^{\vee } )$ are dual to $\ell _n$.
\end{definition}

Here, recall that $\mathfrak{g} ^{\vee } $ is the graded dual given by the internal hom:
\[
	(\mathfrak{g} ^{\vee } )^n = \Hom_{R}(\mathfrak{g}, R[n]) = \Hom_{R}(\mathfrak{g} ^{-n} , R) 
.\]
This algebra is \emph{completed} with respect to the filtation by powers of the ideal generated by $\mathfrak{g} [1]^{\vee } $. This filtration will be important in deformation theory.


Since $C_{\bullet} (\mathfrak{g} )$ is a dg algebra, we want to discuss derivations on it. This is because many Chevalley-Eilenberg cochains in this book will model functions on a space (usually a derived space), and so we want to understand vector fields on this space.

For a free commutative $k$-algebra $\Sym (V^*)$, the derivations are given by $\Sym (V^*)\otimes V$. To see this, let $v\in V$, and consider the evaluation pairing $\left\langle -,- \right\rangle : \Sym ^1(V^*)\otimes V \to k$. We can define an operator $\partial _v =\left\langle -,v \right\rangle : \Sym ^1(V^*) \to k$, and extend it via the Leibniz rule to the full algebra $\Sym (V^*)$. The tensor with $\Sym (V^*)$ arises since, given a derivation $D$ and $f\in \Sym (V^*)$, we can define a derivation $fD$ by evaluating $f$ on $D(v)$. This works since $k\subseteq V^*$ via embedding along scalar multipliation with 1. Since the Chevalley-Eilenberg cochains are indeed a free commutative algebra, this definition turns out to be the homotopically correct answer.

\begin{definition}\label{dfn:10-1-5}
	Let $\mathfrak{g} $ be an $L_\infty $ algebra over $R$. Then the derivations of $\mathfrak{g} $ are given by
	\[
		\operatorname{Der} (\mathfrak{g} )=C^{\bullet} (\mathfrak{g} ,\mathfrak{g} [1]) \coloneqq C^{\bullet} (\mathfrak{g} )\otimes \mathfrak{g} [1]
	.\]
	Here, $\mathfrak{g} [1]$ is a $\mathfrak{g} $-module via a shift of the adjoint action.
\end{definition}

\section{Derived deformation theory}

We work with perturbative field theory, meaning we are interested in how fields behave under infinitesimal changes. There is a mathematical formalism to handle such questions in algebraic geometry, known as \emph{deformation theory}, which studies how geometric properties behave under small perturbations. In fact, a main theme of the book is that the BV formalism provides mechanisms for cleanly expressing many aspects of field theory via deformation theory.

Here, we will briefly introduce derived deformation theory, since we work extensively with derived geometry. We start with formal deformation theory, then move on to the derived theory, and conclude with the role of $L_\infty $ algebras in this theory.

\subsection{The formal neighborhood of a point}

Let $\mathcal{S} \subseteq \Top $ denote some category of spaces. The canonical example for us will be schemes, since we are working algebraically. Yoneda's lemma implies that a space $X\in \mathcal{S} $ can be fully understood using the functor $\mathcal{S} (-,X): \mathcal{S} ^{\mathrm{op}}  \to \Set $. We call $\mathcal{S} (-,X)$  the \emph{functor of points} of $X$, and often write it as $h_X$ to avoid excessive notation. This perspective is quite important in algebraic geometry.

Suppose we want to describe what $X$ looks like near some $p\in X$. The topological answer is to take some open neighborhood of $p$. However, we probably want that open neighborhood to be an object of $\mathcal{S} $ (eg. an open manifold, an open subscheme, etc). Thus, the funtor of points should capture this structure.

We now make this more precise. Fix $\mathcal{S} =\mathcal{S} \mathrm{ch}_{\mathbb{C} } $ -- the category of schemes over $\mathbb{C} $. Here recall that a scheme $X$ over $\mathbb{C} $ means that the space is equipped with a morphism $X \to \operatorname{Spec} \mathbb{C} $. For an affine scheme $X=\Spec R$, this is asking that the underlying ring $R$ be a $\mathbb{C} $-algebra. Thus, $\mathcal{S} \mathrm{ch}_\mathbb{C} =(\mathcal{S} \mathrm{ch} \downarrow\Spec\mathbb{C} )$.

The structure sheaf of any scheme $X$ over $\mathbb{C} $ takes values in commutative $\mathbb{C} $-algebras. Additionally, it's a useful fact that the value of the functor of points $h_X$ is determined by its value on the subategory $\mathcal{A}\mathrm{ff}_{\mathbb{C} } $ of affine $\mathbb{C} $-schemes. Thus, we can view $h_X$ as a functor on $\mathcal{A} \mathrm{ff}_{\mathbb{C} } ^{\mathrm{op}} $ Since an affine scheme yields a $\mathbb{C} $-algebra, there is an equivalence $\mathcal{A} \mathrm{ff}_{\mathbb{C} } \simeq \mathcal{C} \Alg_{\mathbb{C} } ^{\mathrm{op}} $. Therefore, a $\mathbb{C} $-scheme $X$ yields a funtor $h_X : \mathcal{C}\Alg_{\mathbb{C} }  \to \Set $.

Since $\mathbb{C} $ is a field, as a space $\Spec\mathbb{C} \cong *$. Additionally, the structure sheaf is simply
\[
	\mathcal{O}(*) = \mathbb{C} [z\mid V(z)\subseteq V(0)]^{-1} 
.\]
Since $* = V(0)$, we obviously have that $\mathcal{O} (*)\cong \mathbb{C} $. Thus, it makes sense to define a point $p\in X$ of a $\mathbb{C} $-scheme $X$ as a morphism $\Spec\mathbb{C} \to X$ mapping $*\mapsto p$. Since every point has an affine neighborhood, it suffices to understand points of affine schemes. 

We would like to start doing deformation theory by defining a \emph{formal neighborhood of a point}. The idea is as follows: can we ``thicken'' a point by an infinitesimal amount? Preferably, without appealing to open sets? The answer will be given by appealing to a less trivial structure sheaf than that on $\Spec\mathbb{C} $. In particular, we will look for schemes whose space is a point, but whose structure sheaf is not $\mathbb{C} $. We will be most interested in when the structure sheaf looks like an extension of $\mathbb{C} $ by nilpotents, since the intuition for an infinitesimal $dx$ is that $d^2x=0$. This is a much more versatile approach than using neighborhoods, since the Zariski topology is quite coarse.

\begin{definition}\label{dfn:10-2-1}
	A commutative $\mathbb{C} $-algebra $A$ is \emph{Artinian} if $A$ is finite-dimensional as a $\mathbb{C} $-vector space. A \emph{local algebra} $A$ (meaning $A$ has a unqiue maximal ideal $\mathfrak{m} $) is Artinian if and only if there is some $n\geq 0$ such that $\mathfrak{m} ^n=0$.
\end{definition}

We interpret the scheme of an Artinian local algebra as a point with ``infinitesimal directions''. A result in commutative algebra is that for any finite dimensional algebra over a field, all prime ideals are maximal. It then follows that if $(A,\mathfrak{m} )$ is an Artinian local ring, then $\Spec A \cong *$ as a space. Since elements of $\mathfrak{m} $ are ``sufficiently small'' in the sense that $f^n=0$ for all $f\in \mathfrak{m} $ and some $n\in\mathbb{Z} $, we can think of these as sort of infinitesimal operators on the point $*$.

We can once again show that the structure sheaf is simply $\mathcal{O} (*)=A$. Let $\mathcal{A} \mathrm{rt}_{\mathbb{C} } $ denote the category of local Artinian algebras, which we will think of as infinitesimal thickenings of a point.

Every point $p : \Spec\mathbb{C}  \to \Spec R$ corresponds to a map $P : R\to \mathbb{C} $. Since any local Artinian algebra has a distinguished quotient map $Q : A \to A/\mathfrak{m} \cong \mathbb{C} $, postcomposition by $Q$ yields a functor
\[
	h_p : \mathcal{A} \mathrm{rt}_{\mathbb{C} }  \to \Set 
\]
\[
	(A,\mathfrak{m} )\mapsto \left\{ F : R \to A \mid P=Q\circ F \right\} 
.\]
Geometrically, this means $p$ is the composite
\[\begin{tikzcd}[row sep = 2.5em, column sep = 2.5em]
\Spec\mathbb{C} \ar[" \Spec Q ", r]  & \Spec A \ar[" \Spec F ", r]  & \Spec R.
\end{tikzcd}\]
The map $F$ describes some way to infinitesimally extend the point $p$.

Consider the \emph{dual numbers}
\[
	\mathbb{D} \coloneqq \frac{\mathbb{C} [\varepsilon ]}{(\varepsilon ^2)},
\]
which is easily seen to be local Artinian with maximal ideal $(\varepsilon )$. We can show that if $X$ is a scheme, then for any $p\in X$, the set $h_p(\mathbb{D})$ is the tangent space to $p$. Here, $h_p = h_{\Spec\mathbb{C} } $

\begin{proposition}\label{prp:10-2-2}
	Let $P : R \to \mathbb{C} $ be a map of $\mathbb{C} $-algebras, and let $p=\Spec P : \Spec \mathbb{C}  \to \Spec R$ be a point in $\Spec R$. Then
	\[
		h_p(A) \cong \mathcal{C} \Alg_{\mathbb{C} } (\widehat{R}_p,A),
	\]
	\[
		\widehat{R} _p \coloneqq \lim \frac{R}{\mathfrak{m} _p^n} 
	.\]
	Here, $\mathfrak{m} _p = \Ker P$.
\end{proposition}

In other words, the functor $h_p$ is not repesented by a local Artinian algebra, but when passing to the larger category $\mathcal{C} \Alg_{\mathbb{C} } $, we can find a representation as a limit over quotients of $R$. As a regular ring, $\widehat{R}_p$ is given by formal power series, hence motivating further the idea of a formal neighborhood.

There are a few characterizing properties of $h_p$ which are useful in deformation theory, and axiomatizing them will allow us to pass to the derived world. First, note that $h_p(\mathbb{C} )$ is a point by definition: since the quotient $Q$ is the identity, the only function $F$ such that $P=Q\circ F$ is $F=P$. Thus, $h_p(\mathbb{C} )$ is the point $p$.

Next, note that we can study $h_p$ in stages, using a process known as \emph{Artinian induction}. There is a filtration
\[
	A \supseteq \mathfrak{m} \supseteq \mathfrak{m} ^2\supseteq \cdots \supseteq \mathfrak{m} ^n= 0
.\]
Thus, every Artinian algebra can be constructed iteratively by a sequence of small extensions, i.e. a particular short exact sequence.

A third property is that a pullback of local Artinian algebras is Artinian, allowing us to say that
\[
	h_p(A\times_CB)\cong h_p(A)\times_{h_p(C)} h_p(B)
.\]
In other words, $h_p$ preserves pullbacks.

\subsection{Formal moduli spaces}

The functor of points perspective allows us to generalize algebraic geometry by changing the source and target categories of $h_p$. For instance, stacks are functors $\mathcal{C} \Alg_{\mathbb{C} } \to \Grp\dd$, allowing us to capture the notion of symmetries on a scheme. We can enhance further, considering a map from dg commutative $\mathbb{C} $-algebras to $\infty $-groupoids, or more generally, simplicial sets. This is the perspective of derived deformation theory.

\begin{definition}\label{dfn:10-2-3}
	An \emph{Artinian} differential graded commutative $\mathbb{C}$-algebra $A$ (or simply a dg Artinian algebra) is a differential graded commutative $\mathbb{C} $-algebra such that
	\begin{enumerate}
		\item Each degree $A^k$ is finite dimensional as a $\mathbb{C} $-vector space, and there exists $k \ll 0$ sufficiently small such that $\dim_{\mathbb{C} } A^k=0$, and $A^k=0$ for \emph{all} $k>0$.
		\item There is a unique maximal ideal $\mathfrak{m} $ of $A$ which is closed under the differential, and $A / \mathfrak{m} \cong \mathbb{C} $.
	\end{enumerate}
	Morphisms are maps of differential graded commutative $\mathbb{C} $-algebras. The category of differential graded Artinian algebras is denoted $\dg \mathcal{A} \mathrm{rt}_{\mathbb{C}} $.
\end{definition}
Note that locality is part of our definition, since we are most interested in local rings. Additionally, note that we have required $A$ be concentrated in nonpositive degrees. This allows us to appeal to the Dold-Kan correspondence, since we would like to work homotopically, and thus $A$ should be a simplicial commutative algebra.

We can now provide an abstract characterization of the functor of points.

\begin{definition}\label{dfn:10-2-4}
	A \emph{formal moduli problem} is a functor
	\[
		F : \dg \mathcal{A} \mathrm{rt}_{\mathbb{C} }  \to \sSet 
	\]
	such that
	\begin{enumerate}
		\item $F(\mathbb{C} )$ is a contractible Kan complex.
		\item $F$ sends surjections of dg Artinian algebras to fibrations of simplicial sets.
		\item The canonical map $F(A\times _CB)\to F(A)\times _{F(C)} F(B)$ is a weak equivalence.
	\end{enumerate}
\end{definition}

Recall that in the model structure on bounded above cochains, fibrations are exactly the surjections. Thus, $F$ preserves fibrations, and by Ken Brown's lemma $F$ preserves all weak equivalences between fibrant objects. All bounded above cochains are fibrant. Thus, $F$ preserves weak equivalences when it takes values in Kan complexes, but not necessarily for general simplicial sets.

However, in the case of pullbacks, (2) implies that $F(A)\times _{F(C)} F(B)$ is the homotopy pullback as computed in the $\infty $-category $\sSet $.

The largest class of examples that we are interested in are as follows. Let $R=\widehat{\Sym } V$ be a dg commutative $\mathbb{C} $-algebra, whose differential $\dd _R$ is a degree 1 derivation. Then there is a unique maximal ideal $\mathfrak{m} $ generated by $V$. There is a formal moduli problem $h_R$ given by
\[
	h_R(A)_n = \dg \mathcal{C} \Alg_{\mathbb{C}} (R,A\otimes_\mathbb{C} \Omega ^{\bullet} (\left| \Delta ^n \right| ))
.\]
Here, $\Omega ^{\bullet} $ is the de Rham complex, and $\left| \Delta ^n \right| $ is the $n$-simplex. The simplicial operators arise from applying the sheaf $\Omega ^{\bullet} $ to the cosimplicial object $\left| \Delta ^{\bullet}  \right| $ in $\Top $. Then simply recall that Hom is covariant in its se4cond argument. Here, we assume maps of dg commutative algebras are unital. Notice how this formal moduli problem looks similar to the functor $h_p$ induced by the functor of points associated to a point.

\subsection{The role of $L_\infty $ algebras in deformation theory}

There is another source of formal moduli problems, and this approach is the core of the Costello-Gwilliam approach to field theory. This follows the Maurer-Cartan equations $MC(\mathfrak{g} )$ of an $L_\infty $ algebra $\mathfrak{g} $.

\begin{definition}\label{dfn:10-2-5}
	Let $\mathfrak{g} $ be an $L_\infty $ algebra. The \emph{Maurer-Cartan equation} (or MC equation) is
	\[
		\sum_{n=1}^{\infty} \frac{1}{n!} \ell _n(\alpha ^{\otimes n} )=0,
	\]
	for $\alpha \in \mathfrak{g} ^1$. This equation \emph{is not necessarily true for all $\alpha $}.
\end{definition}

Let $\mathfrak{g}$ a Lie algebra, and $M$ a manifold. Then $\Omega ^{\bullet} (M)\otimes \mathfrak{g} $ is a dg Lie algebra with bracket given by (I assume) the commutator on $\Omega ^{\bullet} $ tensored with that of $\mathfrak{g} $. Let $\alpha \in \Omega ^1 \otimes \mathfrak{g} $ be a $\mathfrak{g} $-connection on the trivial principal $G$-bundle (here $\mathfrak{g} =T_1G$). Then $\alpha $ is \emph{flat} if and only if
\[
	\dd \alpha +\frac{1}{2} [\alpha ,\alpha ]=0
.\]
Note that this is exactly the Maurer-Cartan equation of $\alpha $, and is infact its namesake. Thus, flatness equates to the Maurer-Cartan equation holding.

Here's another perspective. A map $\underline{\alpha } : C^{\bullet} (\mathfrak{g} ) \to \mathbb{C} $ of dg algebras is determined by its action on the generators $\mathfrak{g}^{\vee }[-1] =\mathfrak{g} [1]^{\vee } $ of $C^{\bullet} (\mathfrak{g} )$. Hence, $\underline{\alpha } $ is a linear map $\mathfrak{g} ^{\vee } [-1]\to \mathbb{C} $. Equivalently, this is a degree 1 element $\alpha \in \mathfrak{g} $, and $\underline{\alpha } (\beta ) = \beta (\alpha )$. The condition $\underline{\alpha } \circ \dd =0 $ is precisely the Maurer-Cartan equation for $\alpha $.

We can also use the coalgebra $C_{\bullet} (\mathfrak{g} )$. A solution to the MC equation for $\alpha \in \mathfrak{g} $ is equivalent to a map of dg coalgebras $\widetilde{\alpha } : \mathbb{C}  \to C_{\bullet} (\mathfrak{g} )$. This once again simply amounts to unwinding definitions.

We can work even more generally. Let $\mathfrak{g} $ be an $L_\infty $ algebra, and $A$ any dg commutative algebra. Define an $L_\infty $ structure on $\mathfrak{g} \otimes A$ as follows:
\[
	\ell_1(x\otimes a) = \ell _1^{\mathfrak{g} } (x)\otimes a + (-1)^{\left| x \right| } x\otimes \dd_Aa,
\]
\[
	\ell_{n\geq 2} (x_1 \otimes a_1, \ldots , x_n \otimes a_n) = (-1)^{\sum_{1\leq i<j\leq n} \left| a_i \right| \left| x_j \right| } \ell _n^{\mathfrak{g} } (x_1, \ldots ,x_n) \otimes (a_1\cdots a_n)
.\]
Here, $\ell _n^{\mathfrak{g} } $ denotes the bracket on $\mathfrak{g} $. The condition for $\ell _1$ makes it obviously the usual differential on the tensor product, and the higher brackets act only on $\mathfrak{g} $.

Unwinding definitions shows the following statement: a solution $\alpha $ to the MC equation corresponds to a map of dg commutative algebras $\underline{\alpha } : C^{\bullet} (\mathfrak{g}) \to A$ \emph{and} a map of dg cocommutative coalgebras $\widetilde{\alpha } : A \to C_{\bullet} (\mathfrak{g} )$. As maps compose, these elements play sufficiently nicely for us to have a functor.
\begin{definition}\label{dfn:10-2-6}
	Let $\mathfrak{g} $ be an $L_\infty $ algebra. The \emph{Maurer-Cartan functor} of $\mathfrak{g} $ is the functor
	\[
		MC_{\mathfrak{g}}: \dg \mathcal{A} \mathrm{rt}_{\mathbb{C} }  \to \sSet 
	\]
	given by mapping $(A,\mathfrak{m} )$ to the simplicial set whose $n$-simplices are solutions to the Maurer-Cartan equation in $\mathfrak{g} \otimes \mathfrak{m} \otimes \Omega ^{\bullet} (\Delta ^n)$.
\end{definition}

Tensoring with the nilpotent ideal $\mathfrak{m} $ makes $\mathfrak{g}\otimes \mathfrak{m} $ nilpotent. This can be used to show that $MC(\mathfrak{g} )$ takes values in \emph{Kan complexes}, i.e. $MC(\mathfrak{g} )$ factors through the $\infty $-category $\mathcal{S} $ of spaces. There is a result:

\begin{theorem}\label{thm:10-2-7}
	The Maurer-Cartan functor $MC_{\mathfrak{g} } $ is a formal moduli problem.
\end{theorem}

Lurie showed even further that, assuming some nice conditions, every formal moduli problem is weakly equivalent to a Maurer-Cartan functor. Hence, we can use the extremely tractible tools of homological algebra to work with Maurer-Cartan functors and formal moduli problems. We will heavily exploit this throughout the book.

\chapter{Elliptic moduli problems}

\section{Formal moduli problems and Lie algberas}

We now review this theory in the context of field theory. We loosen some of the requirements in order to gain the necessary generality. First, note that we can work over any field $\mathbb{K} $ of characteristic 0, and we adopt this generality here. Next, we modify the definition of a morphism of Artinian algebras to require that the maximal ideal be preserved. Equivalently, a map $(A,m_A) \to (B,m_B)$ is a map of non-unital dg algebras $m_A \to m_B$.

The category $\dg \mathcal{A} \mathrm{rt}_{\mathbb{K} } $ is an $\sSet $-category via the enrichment
\[
	\underline{\dg \mathcal{A} \mathrm{rt} }_{\mathbb{K} } (R,S)_n = \dg \mathcal{C} \Alg_{\mathbb{K} } (m_R,m_S \otimes \Omega ^{\bullet} (\left| \Delta ^n \right| )),
\]
for $\Omega ^{\bullet} (\Delta ^n)$ some commutative algebra model for the cochains on the $n$-simplex. We usually take this to be the de Rham complex when $\mathbb{K} =\mathbb{R} $. We now recast the definition of a formal moduli problem to maintain homotopy coherence.

\begin{definition}\label{dfn:11-1-1}
	A \emph{formal moduli problem} is a $\sSet $-functor $\dg \mathcal{A} \mathrm{rt}_{\mathbb{K} } \to \sSet $ satisfying the hypotheses of \cref{dfn:10-2-4}.
\end{definition}
Since formal moduli problems are now $\sSet $-functors, they can be recast as $\infty $-functors, which are the natural setting given that we are studying cochains up to homotopy. Recall that we may pass to the homotopy coherent nerve to obtain quasi-categories and maps of simplicial sets.

Most authors refer to \cref{dfn:11-1-1} as a \emph{pointed} moduli problem, since $F(\mathbb{K} )$ is contractible (and a Kan complex, but fibrant replacement shows that all simplicial sets have the homotopy type of a Kan complex).

The category of formal moduli problems is simplicially enriched pointwise, and has weak equivalences given pointwise.

\subsection{Formal moduli problems and $L_\infty $ algebras}

The same presentation holds, but to work over a field $\mathbb{K} $, we fix $\Omega ^{\bullet} (\left| \Delta ^{\bullet}  \right| )$ as a model for the cochains on $\left| \Delta ^{\bullet}  \right| $. We also write $B\mathfrak{g} $ for the formal moduli problem associated to the Maurer-Cartan equation for an $L_\infty $ algebra $\mathfrak{g} $.

Since a Maurer-Cartan element $\alpha \in \mathfrak{g} \otimes m_R$ corresponds to a map $C^{\bullet} (\mathfrak{g} )\to R$, we have a way of thinking of $C^{\bullet} (\mathfrak{g} )$ as functions on $B\mathfrak{g} $. Due to some stuff I only vaguely understand, we define
\[
	\Omega ^1(B\mathfrak{g} ) \coloneqq C^{\bullet} (\mathfrak{g} ,\mathfrak{g} ^{\vee } [-1]),\qquad \mathfrak{X} (B\mathfrak{g} )\coloneqq C^{\bullet} (\mathfrak{g} ,\mathfrak{g} [1])
.\]
As noted before, if $\mathfrak{g} $ is finite dimensional, then $\mathfrak{X} $ is the same as derivations of $\mathfrak{g} $, which makese sense. Even if $\mathfrak{g} $ isn't finite dimensional, $\mathfrak{X} $ still controls the deformations of its $L_\infty $ structure up to a shift.

\subsection{The fundamental theorem of deformation theory}

\begin{theorem}\label{thm:11-1-2}
	There is an equivalence of $\infty $-categories between the category of $L_\infty $ algebras and the category of formal moduli problems.
\end{theorem}

This theorem is mostly motivational; the language required to make it precise would take us too far afield into homotopy theory. However, the idea that we can replace any moduli problem with an $L_\infty $ algebra by way of Maurer-Cartan is a profound idea that will motivate most of our discussions.

\subsection{Elliptic moduli problems}

We are interested in moduli problems that describe the solutions to a system of differential equations on a manifold $M$. In particular, these are the Euler-Lagrange equations of motion associated to a local action functional $S$. Such an object yields a sheaf of moduli problems, which can be described by a sheaf of $L_\infty $ algebras.

\begin{definition}\label{dfn:11-1-3}
	Let $M$ be a manifold. A \emph{local $L_\infty $ algebra on $M$} is the following data.
	\begin{enumerate}
		\item A graded vector bundle $L\to M$ with sheaf of smooth sections $\mathcal{L} $.
		\item A differential operator $\dd : \mathcal{L}  \to \mathcal{L} $ of cohomological degree 1, such that $\dd ^2=0$.
		\item A collection of polynomial differential operators
			\[
				\ell _n : \mathcal{L} ^{\otimes n}  \to \mathcal{L} [2-n]
			\]
			for all $n\geq 2$, such that sections of $\mathcal{L} $ are given the structure of an $L_\infty $ algebra.
	\end{enumerate}
\end{definition}

\begin{definition}\label{dfn:11-1-4}
	An \emph{elliptic $L_\infty $ algebra} is a local $L_\infty $ algebra such that $(\mathcal{L} ,\dd )$ is an elliptic complex.
\end{definition}

Given a local $L_\infty $ algebra $\mathcal{L} $, we obtain a presheaf (in fact a homotopy sheaf, but we have no reason to use this fact) $B\mathcal{L} $ of formal moduli problems on $M$, defined by $B\mathcal{L} (U)(R)_n = \mathrm{MC} (\mathfrak{g} \otimes m_R \otimes \Omega ^{\bullet} (\left| \Delta ^{n}  \right| ))$.

\begin{definition}\label{dfn:11-1-5}
	A \emph{formal pointed elliptic moduli problem}, or simply an elliptic moduli problem, is a sheaf of formal moduli problems on a manifold $M$ which is weakly equivalent to $B\mathcal{L} $, for some elliptic $L_\infty $ algebra $\mathcal{L} $ on $M$.
\end{definition}


\section{Examples of elliptic moduli problems related to scalar field theories}

\subsection{The free scalar field theory}

We start by investigating the free scalar field, which we are well-familiar with, before adding interactions. Let $M$ be a Riemannian manifold with metric $g$, and let $\Delta $ be the Laplacian. The text \cite{cg2} differs from \cite{cg1} in the degrees of the fields, but this is fine. The elliptic complex $\mathcal{L} $ is given by
\[
	\mathcal{L} =\left( 
	\begin{tikzcd}
		{C^\infty_M[-1]}\ar[" \Delta  ", r] & {C^\infty _M[-2]}
	\end{tikzcd}
	\right) 
.\]
The brackets $\ell _{n\geq 2} =0$, meaning this is an abelian differential graded Lie $\mathbb{R} $-algebra.

Let's investigate the associated elliptic moduli problem $B\mathcal{L} $. Let $R$ be an ordinary Artinian $\mathbb{R} $-algebra, viewed as a dg Artinian $\mathbb{R} $-algebra by inserting it into degree 0. Let $m$ be the unique maximal ideal. Then
\[
	B\mathcal{L} (U)(R)_n = \mathrm{MC}(\mathcal{L} (U)\otimes m \otimes \Omega ^{\bullet} (\left| \Delta ^n \right| ))
.\]
Since all higher brackets vanish, the Maurer-Cartan equations are simply
\[
	\dd \alpha =0
.\]
This makes the simplicial set incredibly tractible. We start with the 0-simplices. Since $\Delta ^0 \cong *$, we have $\Omega ^{\bullet} (\left| \Delta ^0 \right| ) \cong \mathbb{R} [0]$, so we may drop that term from the tensor. The differential on $\mathcal{L} (U)\otimes _{\mathbb{R} } m$ is simply $\dd (\alpha \otimes r) = \dd \alpha \otimes r$, since $\dd $ acts as a derivation on the tensor product, and thus annihilates any $r\in m$. Additionally, recall that Maurer-Cartan equation is only defined for $\mathfrak{g} ^1$, which is $C^\infty _M[-1]$. Thus,
\[
	B\mathcal{L} (U)(R)_0 \cong \left\{ \alpha \otimes m\in \mathfrak{g} \otimes m\mid \dd \alpha = 0 \right\} \cong \Ker \Delta \otimes_{\mathbb{R} } m \cong H^0(\mathcal{L} (U))\otimes _{\mathbb{R} } m
.\]
We claim that all higher simplicies are constant. Indeed, for $\phi \in \mathcal{L} (U)\otimes_{\mathbb{R} } m\otimes _{\mathbb{R} } \Omega ^{\bullet} (\left| \Delta ^n \right| )$ to have degree 1, then $\phi \in \mathcal{L} (U)^1 \otimes m\otimes \Omega ^0(\left| \Delta ^n \right| )$. For it to be closed, we must have $\dd \phi =0$ and $\dd _{\mathrm{dR} } \phi =0$; the second statement amounts to the statement that the de Rham component of $\phi $ is a constant function. Thus, the $n$-simplicies are simply $\Ker \Delta \otimes _{\mathbb{R} } m\otimes _{\mathbb{R} } \mathbb{R} \cong \Ker \Delta \otimes _{\mathbb{R} } m$.
